\chapter{Regeln --- SWRL}
\label{chap:swrl}

Das vorangehende Kapitel~\ref{chap:owl}, ``\nameref{chap:owl}'' hat gezeigt, wie mittels der Ontologiesprache OWL in der Version 2 eine Ontologie erstellt werden kann. Weiter zeigte das Kapitel wie innerhalb einer Ontologie Klassen, Individuen und Beziehungen aufgebaut werden können. Beim Erstellen solch einer Ontologie, fällt auf, dass der Nutzen der Inferenz mit den beschriebenen Entitäten nicht voll ausgeschöpft werden kann. Damit ein Reasoner die Inferenz vollumfänglich nutzen kann, mussten wir Regeln in Form der Regelsprache \textit{SWRL} einführen.

Das folgende Kapitel basiert, sofern nicht anders vermerkt, auf~\cite{swrl}.

Bei SWRL handelt es sich um eine auf OWL und RuleML~\footnote{www.ruleml.org} basierende Regelsprache. Sie erlaubt es Regeln in Form von OWL-Konzepten auszudrücken und bietet dadurch vielfältige Möglichkeiten der Inferenz.

\section{Aufbau}
\label{sec:swrl_aufbau}
Eine SWRL-Regel besteht aus einem Kopf und einem Körper. Hierbei folgt der Kopf immer auf den Körper. Beide bestehen aus positiven Konjunktionen von Atomen:
\lstset{caption={Beispiel von positiven Konjunktionen von Atomen},captionpos=b}
\lstset{language=XML}
\begin{lstlisting}
    $ atom \wedge atom \cdots \rightarrow atom \wedge atom $
\end{lstlisting}

Informal ausgedrückt heisst das: ``Wenn alle Teile des Körpers wahr sind, dann ist auch der Kopf wahr''. SWRL unterstützt keine negierten Atome oder Disjunktionen.

Ein Atom ist Ausdruck der Form:
\lstset{caption={Beispiel eines Atoms},captionpos=b}
\lstset{language=XML}
\begin{lstlisting}
    $ p(\arg1, \arg2, \cdots argn) $
\end{lstlisting}
$p$ ist hierbei ein Prädikat, $\arg1$, $\arg2$ und $argn$ sind die Argumente bzw. Terme des Ausdrucks.

Prädikate können OWL-Klassen, -Eigenschaften oder -Datentypen sein. Argumente können OWL-Individuen, -Werte oder darauf verweisende Variablen sein.

\noindent\rule[1ex]{\textwidth}{1pt}
\begin{wrapfigure}[4]{l}{0.1\textwidth}
    \vspace{-12pt}
    \includegraphics[width=0.1\textwidth]{bilder/owl.png}
\end{wrapfigure}
In SRWL können Atome ausschliesslich mit UND verknüpft werden. Eine ODER-Verknüpfung gibt es nicht. Die UND-Verknüpfung wird mittels eines Kommas \textit{,} vorgenommen. Eine ODER-Verknüpfung wird duch mehrere Regeln mit unterschiedlichen Regelkörpern, aber gleichen Regelköpfen (Schlussfolgerungen) abgebildet.

\noindent\rule[1ex]{\textwidth}{1pt}

\subsection{Atomare Typen}
\label{subsec:swrl_aufbau_atomaretypen}
SWRL unterscheidet zwischen den folgenden atomaren Typen:
\begin{itemize}
    \item Klassen
    \item Eigenschaften von Individuen
    \item Wertespezifischen Eigenschaften
    \item Unterscheidung von Individuen
    \item Gleichsetzungen von Individuen
    \item Wertebereichen
    \item Vordefinierten Atomen
\end{itemize}

\subsubsection{Klassen}
\label{ssubsec:swrl_aufbau_atomaretypen_klassen}
Ein Klassenatom besteht aus einem Prädikat, sowie nur einem Argument. Das Prädikat ist eine OWL-Klasse, das Argument ein OWL-Individuum oder eine Variable.
\lstset{caption={Beispiele von Klassen-Atomen},captionpos=b}
\lstset{language=XML}
\begin{lstlisting}
    Abenteuerreise(?x)
    Ausflug(Seilpark_Balmberg)
\end{lstlisting}

Möchte man nun mittels SWRL z.B. aussagen, dass alle Abenteuerreisen auch Reisen sind, so geschieht dies mittels:
\lstset{caption={Beispiel der SWRL-Regel, dass alle Abenteuerreisen auch Reisen sind},captionpos=b}
\lstset{language=XML}
\begin{lstlisting}
    Abenteuerreise(?x) $\rightarrow$ Reise(?x)
\end{lstlisting}

\subsubsection{Eigenschaften von Individuen}
\label{ssubsec:swrl_aufbau_atomaretypen_eigenschaftenvonindividuen}
Ein Atom, welches Eigenschaften von Individuen darstellt, besteht aus einem Prädikat und zwei Argumenten. Das Prädikat ist eine OWL-Eigenschaft bzw.\ eine Relation zwischen zwei Individuen. Die Argumente sind somit ein OWL-Individuum oder eine Variable.
\lstset{caption={Beispiele von Atomen zur Darstellung von Eigenschaften von Individuen},captionpos=b}
\lstset{language=XML}
\begin{lstlisting}
    anzahlTeilnehmer(?a, ?anzahl)
    erfordert(?a, "mut")
\end{lstlisting}

Wir definieren als praktisches Beispiel ein Individuum als \textit{Teamevent}, wenn es mindestens 4 maximal aber 20 Teilnehmer zulässt. Weiter muss es über die Eigenschaft der Teambildung verfügen. Drückt man dies per SWRL-Regel aus, ergibt sich folgender Ausdruck:

\lstset{caption={Beispiel der SWRL-Regel um ein Objekt $?a$ als Teamevent zu definieren},captionpos=b}
\lstset{language=XML}
\begin{lstlisting}
    anzahlTeilnehmer(?a, ?anzahl), greaterThanOrEqual(?anzahl, 4), lessThanOrEqual(?anzahl, 20), bietet(?a, "teambildung") $\rightarrow$ teamevent(?a, true)
\end{lstlisting}

\subsubsection{Wertespezifische Eigenschaften}
\label{ssubsec:swrl_aufbau_atomaretypen_wertespezifischeeigenschaften}
Ein Atom, welches wertespezifische Eigenschaften darstellt, besteht aus einem Prädikat sowie zwei Argumenten. Das Prädikat ist eine wertespezifische Eigenschaft. Das erste Argument ein OWL-Individuum, das zweite Argument ist ein OWL-Datenwert.
\lstset{caption={Beispiele von Atomen zur Darstellung von wertespezifischen Eigenschaften von Individuen},captionpos=b}
\lstset{language=XML}
\begin{lstlisting}
    erfordert(?a, "mut")
    durchschnittspreis(?r, ?preis)
\end{lstlisting}

Dass ein Objekt als aufregend (actionreich) gilt, wenn es Mut erfordert, kann mit folgender Regel beschrieben werden:
\lstset{caption={Beispiel einer Regel, welche ausdrückt, dass ein Objekt actionreich ist, wenn es Mut erfordert},captionpos=b}
\lstset{language=XML}
\begin{lstlisting}
    erfordert(?a, "mut") $\rightarrow$ action(?a, true)
\end{lstlisting}

\subsubsection{Unterscheidung von Individuen}
\label{ssubsec:swrl_aufbau_atomaretypen_unterscheidungvonindividuen}
Ein Atom, welches die Unterscheidung von Individuen bezeichnet, besteht aus einem Prädikat sowie zwei Argumenten. Das Prädikat ist das \textit{differentFrom}-Prädikat, die Argumente sind OWL-Individuen.
\lstset{caption={Beispiele von Atomen zur Unterscheidung von zwei Individuen},captionpos=b}
\lstset{language=XML}
\begin{lstlisting}
    differentFrom(?x,?y)
    differentFrom(Seilpark_Balmberg,Seilpark_Pilatus)
\end{lstlisting}
Eine genauere Beschreibung der Verwendung bzw.\ des Sinnes dieses Atoms findet sich unter~\ref{sec:swrl_openworldassumption}.

\subsubsection{Gleichsetzung von Individuen}
\label{ssubsec:swrl_aufbau_atomaretypen_gleichsetzungvonindividuen}
Das Atom, welches die Gleichheit von zwei Individuen kennzeichnet, besteht aus einem Prädikat sowie zwei Argumenten. Das Prädikat ist das \textit{sameAs}-Prädikat, die Argumente sind OWL-Individuen.
\lstset{caption={Beispiele von Atomen zur Gleichsetzung von zwei Individuen},captionpos=b}
\lstset{language=XML}
\begin{lstlisting}
    sameAs(?x,?y)
    sameAs(Seilpark, Seilpark_Balmberg)
\end{lstlisting}

\subsubsection{Wertebereiche}
\label{ssubsec:swrl_aufbau_atomaretypen_wertebereiche}
Das Atom, welches einen Wertebereich kennzeichnet, besteht aus einem Datentyp oder einer Menge von Literalen sowie einem Argument.
\lstset{caption={Beispiele von Atomen zur Kennzeichnung eines Wertes},captionpos=b}
\lstset{language=XML}
\begin{lstlisting}
    xsd:int(?x)
    [3, 4, 5](?x)
\end{lstlisting}
Im ersten Beispiel ist das Objekt, gekennzeichnet durch die Variable $?x$, eine ganzzahlige Zahl.\\
Im zweiten Beispiel hat das Objekt einen der Werte 3, 4 oder 5.

\subsubsection{Vordefinierte Atome}
\label{ssubsec:swrl_aufbau_atomaretypen_vordefinierteatome}
Bei den vordefinierten Atomen handelt es sich um Prädikate, welche ein oder mehrere Argumente entgegennehmen. Sofern alle Argumente dem Prädikat genügen, geben sie den Wahrheitswert ``richtig'' zurück. Andernfalls ``falsch''.\\
Die vordefinierten Atome haben das Präfix $swrldb$. Eine detaillierte Beschreibung aller vordefinierter Atome findet sich unter~\footnote{\url{http://www.daml.org/2004/04/swrl/builtins}}.

Möchte man ausdrücken, dass zum Beispiel Objekte, welche einen Durchschnittspreis zwischen 20 und 30 haben, $preiswert$ (wertespezifische Eigenschaft des Attributes $preissegment$) sind, so geschieht dies mittels der SWRL-Regel:
\lstset{caption={Beispiel einer SWRL-Regel, welche besagt, dass ein Objekt die wertespezifische Eingenschaft \textit{preiswert} hat},captionpos=b}
\lstset{language=XML}
\begin{lstlisting}
    durchschnittspreis(?r, ?preis), lessThanOrEqual(?preis, 30), greaterThan(?preis, 20) $\rightarrow$ preissegment(?r, "preiswert")
\end{lstlisting}

\newpage

\section{Open world assumption}
\label{sec:swrl_openworldassumption}
Wie auch OWL geht SWRL von der so genannten \textit{``open world assumption''} aus. Diese besagt, der Wahrheitswert einer Aussage kann unabhängig des effektiven Wahrheitwertes wahr sein.

Möchte man zum Beispiel ausdrücken, zwei OWL-Individuen sind Kollaborateure (sie arbeiten zusammen an einer Publikation) gilt folgende Regel:

\lstset{caption={Beispiel einer SWRL-Regel zum Ausdrücken von Kollaboration},captionpos=b}
\lstset{language=XML}
\begin{lstlisting}
    $ Publication(?p) , hasAuthor(?p,?y) , hasAuthor(?p,?z) \rightarrow collaboratesWith(?y,?z) $
\end{lstlisting}

Aufgrund  der ``open world assumption'' von OWL ist es nicht möglich auszusagen, dass zwei OWL-Individuen automatisch unterschiedlich sind, wenn sie verschiedene Namen haben. Mittels den \textit{sameAs}- und \textit{differentFrom}-Relationen bietet SWRL die Lösung des Problemes:

\lstset{caption={Beispiel einer SWRL-Regel zum Ausdrücken von Kollaboration mit expliziter Unterscheidung zwischen den Kollaborateuren},captionpos=b}
\lstset{language=XML}
\begin{lstlisting}
    $ Publication(?p) , hasAuthor(?p, ?y) , hasAuthor(?p, ?z) , differentFrom(?y, ?z) \rightarrow collaboratesWith(?y, ?z) $
\end{lstlisting}

Analog gilt dies für die Gleichsetzung von zwei Individuen. Diese müssen explizit mit der \textit{sameAs}-Relation gleichgesetzt werden.

Aufzählungen können beispielsweise nicht ohne Umstände ausgedrückt werden: Es ist nicht möglich zu sagen, ob eine Publikation nur von einem einzigen Autor stammt. Es kann durchaus sein, dass nur ein Autor der Publikation bekannt ist, diese aber noch über weitere Autoren verfügt.

Dies bezeichnet den Kern der ``open world assumption''. Der effektive Wahrheitswert ist beispielsweise, die Publikation $?p$  hat mehrere Autoren ($?x, ?y$ und $?z$). Die Aussage, dass die Publikation $?p$ einen bestimmten Autor $?z$ hat, ist dennoch wahr, auch wenn die übrigen Autoren noch nicht bekannt sind.

\noindent\rule[1ex]{\textwidth}{1pt}
\begin{wrapfigure}[5]{l}{0.1\textwidth}
    \vspace{-14pt}
    \includegraphics[width=0.1\textwidth]{bilder/elephant.png}
\end{wrapfigure}
In diesem Kapitel wurden \textit{Regeln} eingeführt, die es ermöglichen die im~\autoref{chap:graph_data} (\nameref{chap:graph_data}) gefundenen Kriterien \textit{familienfreundlich}, \textit{regional} und \textit{actionreich} als Schlussfolgerungen zu definieren. Vorgängig müssen diese jedoch als Attribute bzw. DataProperties (sepzifischer Datenwert eines Objektes) definiert werden.

Wir haben die Definition wie folgt gewählt:

Ein Objekt muss über das Attribut \textit{alter} verfügen und der Wert dieses Attributes muss grösser oder gleich der Zahl zwei sein. Ist dies gegeben, erhält ein Objekt automatisch das Attribut \textit{familienfreundlich} mit dem Wahrheitswert (Boolean) \textit{wahr} bzw. \textit{true}.

\begin{lstlisting}[caption={Beispiel der SWRL-Regel für die Eigenschaft \textit{familienfreundlich}},captionpos=b,language=XML]
    alter(?a, ?alter), swrldb:greaterThanOrEqual(?alter, 2) $\Rightarrow$ familienfreundlich(?a, true)
\end{lstlisting}

Ein Objekt gilt dann als \textit{regional}, wenn es über einen Standort verfügt und die Region dieses Standortes 50 oder weniger Kilometer vom Ausgangspunkt (welcher dynamisch bestimmt sein kann, also abhängig vom Abfragesteller) entfernt ist.

\begin{lstlisting}[caption={Beispiel der SWRL-Regel für die Eigenschaft \textit{regional}},captionpos=b,language=xml]
    hatStandort(?ausflug, ?ort), hatRegion(?land, ?region), hatOrt(?region, ?ort), distanzZuAusgangspunkt(?region, ?distanz), swrldb:lessThanOrEqual(?distanz, 50) $\Rightarrow$ regional(?ausflug, true)
\end{lstlisting}

Eine Regel kann in SWRL keine ODER-Verknüpfung enthalten, wie in~\autoref{sec:swrl_aufbau} (\nameref{sec:swrl_aufbau}) beschrieben. Die ODER-Verknüpfungen bildet man ab, durch mehrfache Definition von Regeln für die selbe Schlussfolgerung. Der Regelkopf (das heisst die Schlussfolgerung) muss immer gleich sein, die Bedingungen (das heisst die Regelköper) dürfen verschieden sein.

Wir wenden an: Ein Objekt erhält das Attribut \textit{actionreich}, wenn es Mut oder Geschicklichkeit erfordert oder wenn es Nervenkitzel verursacht. Die Eigenschaft \textit{nervenkitzel} ist ihrerseits eine Folgerung aus der Erfordernis von Mut.

Es findet also eine transitive Ableitung statt: $A \sim B \wedge B \sim C \Rightarrow A \sim C$ bzw. $erfordertMut \Rightarrow birgtNervenkitzel \wedge birgtNervenkitzel \Rightarrow istActionreich$ somit gilt $erfordertMut \Rightarrow istActionreich$.

\begin{lstlisting}[caption={Beispiel der SWRL-Regel für die Eigenschaft \textit{actionreich}},captionpos=b,language=XML]
    erfordert(?a, "mut") $\Rightarrow$ actionreich(?a, true)
    nervenkitzel(?a, true) $\Rightarrow$ actionreich(?a, true)
    erfordert(?a, "geschicklichkeit") $\Rightarrow$ actionreich(?a, true)
\end{lstlisting}

Es sind jetzt alle Komponenten vorhanden, um Ausflüge abbilden zu können, die den in~\autoref{chap:graph_data} (\nameref{chap:graph_data}) genannten Kriterien entsprechen.

Um nach Ausflügen mit den gewählten Kriterien suchen zu können, benötigt man jedoch eine bestimmte Abfragesprache. Diese wird im kommenden Kapitel aufgezeigt.

\noindent\rule[1ex]{\textwidth}{1pt}
